\documentclass[10pt,a4paper]{amsart}
\usepackage[margin=19mm]{geometry}
\usepackage{amsmath,amssymb,mathtools}
\usepackage[scr=rsfs,cal=euler]{mathalfa}
\usepackage{xfrac}
\usepackage[unicode,hidelinks,bookmarks=false]{hyperref}
\newcommand{\ie}{i.e.\ }
\newcommand{\cf}{cf.\ }
\newcommand{\resp}{respectively\ }
\newcommand{\Subsection}{\subsection}
\providecommand{\nobreakdash}{\nobreak\hbox{-}}
\setlength{\emergencystretch}{2em}
\multlinegap0pt
\usepackage{mathtools,array}
\hypersetup{hypertex=true,
	colorlinks=true,
	linkcolor=blue,
	anchorcolor=blue,
	citecolor=blue}
\usepackage{bm}
\allowdisplaybreaks[4]
\def\mo{{\rm{mod}\ }}
\def\lcm{\rm{lcm}}
\def\ord{{\rm ord}}
\def\sgn{{\rm sgn}}
\def\gen{{\rm gen}}
\def\spn{{\rm spn}}
\def\Aut{{\rm Aut}}
\def\cs{\ldots}
\def\supp {\rm supp}
\def\colon{{:}\;}
\def\Proof{\noindent{\it Proof}}
\def\Z{\mathbb Z}
\def\C{\mathbb C}
\def\N{\mathbb N}
\def\Q{\mathbb Q}
\def\R{\mathbb R}
\def\E{\mathbb E}
\def\cP{\mathcal P}
\def\O{{\mathcal{O}}}
\def\rank{{\rm rank}}
\def\Gal{{\rm Gal}}
\def\Tr{{\rm Tr}}
\def\Norm{{\rm Norm}}
\def\char{{\rm char}}
\def\Ker{{\rm Ker}}
\def\a{{\bm a}}
\def\u{{\bm u}}
\def\v{{\bm v}}
\def\w{{\bm w}}
\def\x{{\bm x}}
\def\y{{\bm y}}
\def\e{{\bm e}}
\def\0{{\bm 0}}
\def\i{{\bm i}}
\def\1{{\bf 1}}
\def\alg{{\rm alg}}
\def\div{{\rm div}}
\def\diag{{\rm diag}}
\def\sgn{{\rm sign}}
\def\Re{{\rm Re}}
\def\B{{\rm B}}
\def\Hom{{\rm Hom}}
\def\char{{\rm char}}
\def\id{{\rm id}}
\def\Im{{\rm Im}}
\def\supp{{\rm supp}}
\def\bad{{\rm bad}}
\def\good{{\rm good}}
\def\bchi{{\bm \chi}}
\def\lcm{{\rm lcm}}
\def\diam{{\rm diam}}
\def\pmod #1{\ ({\rm{mod}}\ #1)}
\def\mod #1{\ {\rm mod}\ #1}
\def\Ack{\medskip\noindent {\bf Acknowledgments}}
\def\Del{\medskip\noindent{\bf Declarations}}
\newtheorem{theorem}{Theorem}[section]
\newtheorem{lemma}{Lemma}
\newtheorem{corollary}{Corollary}
\newtheorem{conjecture}{Conjecture}
\newtheorem{proposition}{Proposition}
\newtheorem*{Theorem}{Theorem}
\newtheorem{definition}{Definition}
\newtheorem*{Acks}{Acknowledgments}
\newtheorem{remark}{Remark}
\renewcommand{\thetheorem}{\arabic{section}.\arabic{theorem}}
\renewcommand{\thelemma}{\arabic{section}.\arabic{lemma}}
\renewcommand{\theconjecture}{\arabic{section}.\arabic{conjecture}}
\renewcommand{\theproposition}{\arabic{section}.\arabic{proposition}}
\renewcommand{\theremark}{\arabic{section}.\arabic{remark}}
\renewcommand{\thecorollary}{\arabic{section}.\arabic{corollary}}
\renewcommand{\thedefinition}{\arabic{section}.\arabic{definition}}
\newcommand{\sign}[1]{\mathrm{sign}(#1)}
\begin{document}
\title[Arithmetic progressions in dense subsets of primitive elemen]{Arithmetic progressions in dense subsets of primitive elements of finite fields}
\author{Hai-Liang Wu}
\date{}
\maketitle
\noindent\emph{Source and licence.} Arithmetic progressions in dense subsets of primitive elements of finite fields. Version arXiv:2609.29653v1 (CC BY 4.0). This material is distributed under the Creative Commons Attribution 4.0 International licence, http://creativecommons.org/licenses/by/4.0/, as recorded in the arXiv OAI licence metadata for this exact version and shown on the abstract page https://arxiv.org/abs/2609.29653v1. Full rights notice: licenses/arxiv-2609.29653-CC-BY-4.0.txt.\par\smallskip
\noindent\textit{Editorial scope.} Counted unit: the original \texttt{lemma} environment \texttt{Lem. estimates for nontrival character sums} at source lines 267--273 together with its complete proof at lines 275--396; the delivered document keeps source lines 167--397 so that every referenced label and every citation resolves inside it. Native editable source is the paper's own concatenated TeX; the AMS wrapper is editorial formatting only. No publisher class, upstream script or unrelated artwork is redistributed.
	\begin{theorem}\label{Thm. A}
		Let $k\ge3$ be an integer and $0<\delta\le 1$ be a real number. Then for every sufficiently large prime $p$, any subset $A\subseteq\mathcal{P}_p$ with $\#A/\#\mathcal{P}_p\ge \delta$ contains a non-trivial arithmetic progression of length $k$ over $\mathbb{F}_p$. 
	\end{theorem}
	
	Let 
	$$\Omega_p=\left\{g\in\mathbb{Z}: 1\le g\le p-1\ \text{and $g$ is a primitive root modulo $p$}\right\}$$
	be the lift of $\mathcal{P}_p$ to $\mathbb{Z}$. As a direct consequence of Theorem \ref{Thm. A}, we have the following result.
	
	\begin{corollary}\label{Corollary of Thm. A}
		Let $k\ge3$ be an integer and $0<\delta\le 1$ be a real number. Then for every sufficiently large prime $p$, each subset $B\subseteq\Omega_p$ with $\#B/\#\Omega_p\ge\delta$ contains a nontrivial $k$-term arithmetic progression of length $k$ over $\mathbb{Z}$. 
	\end{corollary}
		
		
		
		Roughly speaking, Theorem \ref{Thm. A} illustrates that every
		subset $A\subseteq\mathcal P_p$ with 
		$$\#A\asymp (\#\mathcal{P}_p)$$ 
		contains a nontrivial arithmetic progression of length $k$ for all sufficiently large $p$. The next theorem will show that Theorem \ref{Thm. A} is optimal at the level of the power of $\#\mathcal{P}_p=\varphi(p-1)$:  the condition $\#A\asymp (\#\mathcal{P}_p)$ cannot be replaced by $\#A\asymp (\#\mathcal{P}_p)^{1-\varepsilon}$ for any fixed real number $\varepsilon\in(0,1)$. 
	
	
	\begin{theorem}\label{Thm. B}
		Fix an integer $k\ge3$ and a real number $\varepsilon\in(0,1)$. For every sufficiently large prime $p$, there is a subset $A\subseteq\mathcal{P}_p$ with $\#A=\lfloor (\#\mathcal{P}_p)^{1-\varepsilon}\rfloor$ such that $A$ contains no nontrivial arithmetic progression of length $k$.
	\end{theorem}
	
	
	\subsection{Outline of the paper} In Section 2, we will briefly introduce the relative Szemer\'edi theorem established by Conlon, Fox and Zhao. The proofs of our main results will be given in Section 3--5 respectively. 
	
	\section{The relative Szemer\'edi theorem due to Conlon, Fox and Zhao}
	\setcounter{lemma}{0}
	\setcounter{theorem}{0}
	\setcounter{equation}{0}
	\setcounter{conjecture}{0}
	\setcounter{remark}{0}
	\setcounter{corollary}{0}
	
	Let $N$ be a sufficiently large integer. For any function $f(x_1,x_2,\cdots, x_m): \left(\mathbb{Z}/N\mathbb{Z}\right)^m\rightarrow\mathbb{R}$ and any nonempty subsets $A_1,\cdots,A_m\subseteq \mathbb{Z}/N\mathbb{Z}$, let 
	$$\mathbb{E}\left[f(x_1,\cdots,x_m): x_1\in A_1,\cdots, x_m\in A_m\right]$$
	be the expectation of $f(x_1,\cdots,x_m)$ when each $x_i$ is chosen uniformly and independently at random from $A_i$. 
	
	For any positive integer $k$, the symbol $[1,k]$ denotes the set $\{1,2,\cdots, k\}$. Conlon, Fox and Zhao  \cite[Definition 2.2]{CFZ} introduce the following $k$-linear forms condition.
	
	\begin{definition}[$k$-linear forms condition]\label{Def. linear forms condition}
		Let $k$ be a positive integer and let $N$ be a sufficiently large integer. A nonnegative function $v_N: \mathbb{Z}/N\mathbb{Z}\rightarrow\mathbb{R}_{\ge 0}$ is said to obey the $k$-linear forms condition if one has 
		$$\mathbb{E}\left[\prod_{j=1}^{k}\prod_{\substack{{\bm \omega}=(\omega_i)\\ i\in[1,k]\setminus\{j\} \\ \omega_i\in \{0,1\}}}v_N\left(\sum_{i\in[1,k]\setminus\{j\}}(i-j)x_i^{(\omega_i)}\right)^{n_{j,{\bm \omega}}}: x_1^{(0)},x_1^{(1)},\cdots, x_k^{(0)},x_k^{(1)}\in\mathbb{Z}/N\mathbb{Z}\right]=1+o(1),$$
		for any choices of exponents $n_{j,{\bm \omega}}\in\{0,1\}$. 
	\end{definition}
	
	Conlon, Fox and Zhao \cite[Theorem 2.4]{CFZ} established the following result. 
	
	\begin{theorem}[Conlon, Fox and Zhao]\label{Thm. CFZ}
		For every integer $k\ge3$ and real number $\delta>0$, there exists $c>0$ such that if $v_N: \mathbb{Z}/N\mathbb{Z}\rightarrow\mathbb{R}_{\ge 0}$ satisfies the $k$-linear forms condition, $N$ is sufficiently large, and $f: \mathbb{Z}/N\mathbb{Z}\rightarrow\mathbb{R}_{\ge 0}$ satisfies $0\le f(x)\le v_N(x)$ for all $x\in \mathbb{Z}/N\mathbb{Z}$ and $\mathbb{E}[f(x): x\in \mathbb{Z}/N\mathbb{Z}]\ge \delta$, then 
		$$\mathbb{E}\left[f(x)f(x+d)\cdots f(x+(k-1)d): x, d\in\mathbb{Z}/N\mathbb{Z}\right]\ge c.$$
	\end{theorem}
	
	
	
	
	
	\section{Proofs of Theorem \ref{Thm. A}}
	\setcounter{lemma}{0}
	\setcounter{theorem}{0}
	\setcounter{equation}{0}
	\setcounter{conjecture}{0}
	\setcounter{remark}{0}
	\setcounter{corollary}{0}
	
	Recall that $\mathcal{P}_p$ is the set of all primitive elements of $\mathbb{F}_p$. We begin with the following known result on the characteristic function of $\mathcal{P}_p$ (cf. \cite{Cohen15,Cohen21}).
	
	\begin{lemma}\label{Lem. characteristic function of P}
		Let $p$ be a prime. Then 
		$$1_{\mathcal{P}_p}(x)=\theta_{p-1}\sum_{d\mid p-1}\frac{\mu(d)}{\varphi(d)}\sum_{\substack{\chi\in\widehat{\mathbb{F}_p^*} \\ \ord(\chi)=d}} \chi(x)=\theta_{p-1}\sum_{\chi\in\widehat{\mathbb{F}_p^*}}c_{\chi}\cdot \chi(x)
		=\begin{cases}
			1  & \mbox{if}\ x\in\mathcal{P}_p,\\
			0 & \mbox{otherwise},
		\end{cases}$$
		where $\mu$ is the M\"obius function, $\varphi$ is the Euler totient function, $\theta_{p-1}=\frac{\varphi(p-1)}{(p-1)}$, and $c_{\chi}=\frac{\mu(\ord(\chi))}{\varphi(\ord(\chi))}$ for any $\chi\in\widehat{\mathbb{F}_p^*}$. 
	\end{lemma}
	
	We now define the function 
	\begin{equation}\label{Eq. definition of vp}
		v_p=\frac{1}{\theta_{p-1}}\cdot 1_{\mathcal{P}_p}: \mathbb{F}_p\rightarrow \mathbb{R}_{\ge 0}.
	\end{equation}
	
	Our next task is to show that $v_p$ satisfies the $k$-linear forms condition in Definition \ref{Def. linear forms condition}. To accomplish this task, we first introduce the well-known Weil theorem (cf. \cite[Theorem 5.41]{LN}).
	
	\begin{lemma}\label{Lem. the Weil Bound}
		Let $\chi\in\widehat{\mathbb{F}_p^*}$ with $\ord(\chi)=d>1$, and let $f(t)\in\mathbb{F}_p[t]$ be a monic polynomial with $f(t)\neq g(t)^d$ for any $g(t)\in\mathbb{F}_p[t]$. Then, for any $a\in\mathbb{F}_p$ we have 
		$$\left|\sum_{x\in\mathbb{F}_p}\chi(af(x))\right|\le (r-1)\sqrt{p},$$
		where $r$ is the number of distinct roots of $f(t)$ in an algebraic closure $\overline{\mathbb{F}_p}$ of $\mathbb{F}_p$.
	\end{lemma}
	
	Alon \cite[Theorem 1.2]{A} obtained the following useful result, which is known as the Combinatorial Nullstellensatz.
	
	\begin{lemma}\label{Lem. Alon}
		Let $F$ be an arbitrary field, and let $f(x_1,\cdots,x_m)\in F[x_1,\cdots,x_m]$. Suppose $\deg(f)=t_1+t_2+\cdots+t_m$, where each $t_i$ is a nonnegative integer, and suppose that the coefficient of $x_1^{t_1}\cdots x_m^{t_m}$ in $f$ is nonzero. Then, if $S_1,\cdots,S_m$ are subsets of $F$ with $\#S_i>t_i$, there are $s_1\in S_1, s_2\in S_2,\cdots, s_m\in S_m$ so that 
		$$f(s_1,s_2,\cdots,s_m)\neq 0.$$
	\end{lemma}
	
	Using Lemmas \ref{Lem. the Weil Bound}--\ref{Lem. Alon}, we can obtain the following result, which will play an important role in verifying that $v_p$ satisfies the $k$-linear forms condition.
	

	\begin{lemma}\label{Lem. estimates for nontrival character sums}
		Let $m,s\ge2$ be integers, and let $L_i(x_1,x_2,\cdots,x_m)\in\mathbb{Z}[x_1,x_2,\cdots,x_m]$ with $\deg(L_i)=1$ for any $1\le i\le s$. Suppose  that $L_i$ and $L_j$ are not proportional over $\mathbb{Q}$ for any $1\le i\neq j\le s$. Then, for every sufficiently large prime $p$ we have 
		\begin{equation}\label{Eq. inequality in Lemma estimates for nontrival character sums}
			\left|\sum_{\x\in\mathbb{F}_p^m}\prod_{i=1}^{s}\chi_i(L_i(\x))\right|\le (s-1)\cdot p^{m-\frac{1}{2}}+\frac{s(s-1)}{2}\cdot p^{m-1},
		\end{equation}
		where $\chi_1,\chi_2,\cdots,\chi_s\in\widehat{\mathbb{F}_p^*}$ are not all trivial. 
	\end{lemma}

	\begin{proof}
		For each integer $1\le i\le s$, let 
		$$L_i(\x)=a_{i1}x_1+a_{i2}x_2+\cdots+a_{im}x_m+b_i=\a_i\cdot \x+b_i,$$
		where $\a_i=(a_{i1},a_{i2},\cdots,a_{im})\neq \0$, and let $\widetilde{\a_i}=(a_{i1},a_{i2},\cdots,a_{im},b_i)$. 
		
		For any prime $p$, we may view $L_1, L_2, \cdots,L_s$ as polynomials over $\mathbb{F}_p$. We first show that for sufficiently large $p$, when viewed as polynomials over $\mathbb{F}_p$, we still have $\deg(L_i)=1$ for any $1\le i\le s$ and $L_i$ and $L_j$ are not proportional over $\mathbb{F}_p$ for any $1\le i\neq j\le s$. In fact, given any $1\le i\neq j\le s$, since $L_i$ and $L_j$ are not proportional over $\mathbb{Q}$, the vectors $\widetilde{\a_i}$ and $\widetilde{\a_j}$ are linearly independent over $\mathbb{Q}$. From this, the rank of the matrix $\begin{bmatrix}
			\widetilde{\a_i}\\
			\widetilde{\a_j}
		\end{bmatrix}$ is equal to $2$, and hence this matrix
		has a nonzero $2\times 2$ minor $D_{ij}$. Applying this and noting that 
		$$d_i=\max\left\{|a_{ij}|: 1\le j\le m\right\}>0$$
		for each $1\le i\le s$, 
		 when 
		$$p>p_0=\max\left\{s+1, \prod_{1\le i<j\le s}|D_{ij}|\cdot \prod_{1\le i\le s}d_i\right\},$$
		it is clear that $\deg(L_i)=1$ for any $1\le i\le s$ and $L_i$ and $L_j$ are not proportional over $\mathbb{F}_p$ for any $1\le i\neq j\le s$.
		
		We next show that for any $p>p_0$, there exists a vector $\v\in\mathbb{F}_p^m$ such that $\v\cdot \a_i\neq 0$ for any $1\le i\le s$. In fact, since $p>p_0$, the vectors $\a_i\in\mathbb{F}_p^m\setminus\{\0\}$ for each $1\le i\le s$. Thus, 
		$$\left\{\v\in\mathbb{F}_p^m: \v\cdot \a_i=0\right\}$$
		is a subspace of $\mathbb{F}_p^m$ with dimension $m-1$ for any $1\le i\le s$. From this and noting that $p>p_0>s$, one can verify that 
		\begin{align*}
			\#\left\{\v\in\mathbb{F}_p^m: \v\cdot \a_i=0\ \text{for some $1\le i\le s$}\right\}
&=\#\bigcup_{i=1}^s\left\{\v\in\mathbb{F}_p^m: \v\cdot \a_i=0\right\}\\
&\le \sum_{i=1}^s\#\left\{\v\in\mathbb{F}_p^m: \v\cdot \a_i=0\right\}\\
&=sp^{m-1}\\
&<p^m.
		\end{align*}
		Hence there exists a vector $\v\in\mathbb{F}_p^m$ such that $\v\cdot \a_i\neq 0$ for any $1\le i\le s$ whenever $p>p_0$. 
		
		From now on, we assume $p>p_0$. Since $m\ge2$, there is a subspace $W$ of $\mathbb{F}_p^m$ with $\dim(W)=m-1$ such that 
		$$\mathbb{F}_p^m=W\bigoplus\mathbb{F}_p\v.$$
		This implies that every vector $\x\in\mathbb{F}_p^m$ can be uniquely written as 
		$$\x=\w+t\v \quad (\w\in W, t\in\mathbb{F}_p).$$
		Hence, for any $\x=\w+t\v\in\mathbb{F}_p^m$ with $\w\in W$ and $t\in\mathbb{F}_p$ and any integer $1\le i\le s$, letting 
		\begin{equation}\label{Eq. definition of ci}
			c_i=\v\cdot\a_i\neq 0,
		\end{equation}
		one can verify that 
		\begin{align}\label{Eq. expression of Li}
			L_i(\x)=L_i(\w+t\v)
&=\a_i\cdot (\w+t\v)+b_i\notag\\
&=L_i(\w)+c_it\notag\\
&=c_i(t-\alpha_i(\w)),
		\end{align}
where 
      \begin{equation}\label{Eq. definition of alpha i}
	\alpha_i(\w)=-\frac{L_i(\w)}{c_i}.
       \end{equation}


Define 
$$W_{\bad}=\left\{\w\in W: \alpha_i(\w)=\alpha_j(\w)\ \text{for some $1\le i\neq j\le s$}\right\}$$
and let $W_{\good}=W\setminus W_{\bad}$. We next obtain an upper bound for $\# W_{\bad}$. For any $1\le i\neq j\le s$, let the polynomial
$$L_{ij}=c_jL_i-c_iL_j.$$
Since $p>p_0$, by the above discussions, the polynomials $L_i$ and $L_j$ are not proportional over $\mathbb{F}_p$. From this and recalling that  $c_i,c_j\neq 0$ are defined by (\ref{Eq. definition of ci}), we immediately obtain $L_{ij}\neq 0$. On the other hand, for any $\x=\w+t\v\in\mathbb{F}_p^m$ with $\w\in W$ and $t\in\mathbb{F}_p$, applying (\ref{Eq. expression of Li}) and (\ref{Eq. definition of alpha i}) one can verify that 
\begin{align}\label{Eq. Lij is independent of tv}
	L_{ij}(\x)=L_{ij}(\w+t\v)
&=c_jL_i(\w+t\v)-c_iL_j(\w+t\v)\notag\\
&=c_j\left(c_i(t-\alpha_i(\w))\right)-c_i\left(c_j(t-\alpha_j(\w))\right)\notag\\
&=-c_jc_i\alpha_i(\w)+c_ic_j\alpha_j(\w)\notag\\
&=c_jL_i(\w)-c_iL_j(\w)\notag\\
&=L_{ij}(\w).
\end{align}
We now claim that there is at least one $\w\in W$ such that $L_{ij}(\w)\neq 0$. Suppose, to the contrary, that $L_{ij}(\w)=0$ for any $\w\in W$. Then (\ref{Eq. Lij is independent of tv}) implies that $L_{ij}(\x)=0$ for any $\x\in\mathbb{F}_p^m$. Thus, applying Lemma \ref{Lem. Alon} with $F=\mathbb{F}_p$ and $S_1=\cdots=S_m=\mathbb{F}_p$, we obtain $L_{ij}=0$, which contradicts the fact that $L_{ij}\neq 0$. By (\ref{Eq. Lij is independent of tv}) and the above discussions, 
$$\left\{\w\in W: L_{ij}(\w)=0\right\}\subsetneqq W$$
is either $\emptyset$ or a proper affine hyperplane in $W$. From this and noting that 
$$\left\{\w\in W: \alpha_i(\w)=\alpha_j(\w)\right\}=\left\{\w\in W: L_{ij}(\w)=0\right\},$$
we obtain 
$$\#\left\{\w\in W: \alpha_i(\w)=\alpha_j(\w)\right\}\le p^{\dim(W)-1}=p^{m-2}.$$
Hence 
\begin{align}\label{Eq. upper bound for W bad}
	\#W_{\bad}
&=\#\left\{\w\in W: \alpha_i(\w)=\alpha_j(\w)\ \text{for some $1\le i\neq j\le s$}\right\}\notag\\
&=\#\bigcup_{1\le i<j\le s}\left\{\w\in W: \alpha_i(\w)=\alpha_j(\w)\right\}\notag\\
&\le \sum_{1\le i<j\le s}\#\left\{\w\in W: \alpha_i(\w)=\alpha_j(\w)\right\}\notag\\
&\le \binom{s}{2}p^{m-2}.
\end{align}

Finally, we prove the inequality (\ref{Eq. inequality in Lemma estimates for nontrival character sums}). As $\chi_1,\cdots,\chi_s$ are not all trivial, we have 
$$l=\lcm\{\ord(\chi_i): 1\le i\le s\}>1.$$
Fix a character $\psi\in\widehat{\mathbb{F}_p^*}$ with $\ord(\psi)=l$. Then there are integers $1\le e_1,\cdots, e_s\le l$ such that $\chi_i=\psi^{e_i}$ for any $1\le i\le s$. Since $\chi_1,\cdots,\chi_s$ are not all trivial, at least one of $e_1,\cdots,e_s$ is less than $l$. By this and (\ref{Eq. expression of Li}), one can verify that 
\begin{align}\label{Eq. sum is equal to good plus bad}
	\sum_{\x\in\mathbb{F}_p^m}\prod_{i=1}^{s}\chi_i(L_i(\x))
&=\sum_{\w\in W}\sum_{t\in\mathbb{F}_p}\prod_{i=1}^{s}\chi_i(L_i(\w+t\v))\notag\\
&=\sum_{\w\in W}\sum_{t\in\mathbb{F}_p}\prod_{i=1}^{s}\chi_i(c_i(t-\alpha_i(\w)))\notag\\
&=\sum_{\w\in W}\sum_{t\in\mathbb{F}_p}\prod_{i=1}^{s}\psi^{e_i}(c_i(t-\alpha_i(\w)))\notag\\
&=\sum_{\w\in W_{\good}}\sum_{t\in\mathbb{F}_p}\prod_{i=1}^{s}\psi^{e_i}(c_i(t-\alpha_i(\w)))+\sum_{\w\in W_{\bad}}\sum_{t\in\mathbb{F}_p}\prod_{i=1}^{s}\psi^{e_i}(c_i(t-\alpha_i(\w))).
\end{align}
	
	For any $\w\in W$, let 
	$$f_{\w}(x)=\prod_{1\le i\le s}\left(x-\alpha_i(\w)\right)^{e_i}\in\mathbb{F}_p[x].$$
	If $\w\in W_{\good}$, then $\alpha_1(\w),\alpha_2(\w),\cdots, \alpha_s(\w)$ are pairwise distinct. Noting that $1\le e_1,\cdots,e_s\le l$ and $\min\{e_i: 1\le i\le s\}<l$, we have 
		$$f_{\w}(x)\neq g(x)^l$$
		for any $g(x)\in\mathbb{F}_p[x]$.  Applying Lemma \ref{Lem. the Weil Bound}, for each $\w\in W_{\good}$ we obtain 
$$\left|\sum_{t\in\mathbb{F}_p}\prod_{i=1}^{s}\psi^{e_i}(c_i(t-\alpha_i(\w)))\right|=\left|\sum_{t\in\mathbb{F}_p}\psi\left(cf_{\w}(t)\right)\right|\le (s-1)\sqrt{p},$$
where 
$$c=\prod_{1\le i\le s}c_i^{e_i}\in\mathbb{F}_p^*.$$
Hence, 
\begin{align}\label{Eq. estimate for good w}
	     \left|\sum_{\w\in W_{\good}}\sum_{t\in\mathbb{F}_p}\prod_{i=1}^{s}\psi^{e_i}(c_i(t-\alpha_i(\w)))\right|
&\le  \sum_{\w\in W_{\good}}\left|\sum_{t\in\mathbb{F}_p}\prod_{i=1}^{s}\psi^{e_i}(c_i(t-\alpha_i(\w)))\right|\notag\\
&\le \#W_{\good}\cdot (s-1)\cdot \sqrt{p}\notag\\
&\le \#W\cdot (s-1)\cdot \sqrt{p}\notag\\
&=(s-1)\cdot p^{m-\frac{1}{2}}.
\end{align}		
On the other hand, by (\ref{Eq. upper bound for W bad}) one can verify that 
\begin{align}\label{Eq. estimate for bad w}
	\left|\sum_{\w\in W_{\bad}}\sum_{t\in\mathbb{F}_p}\prod_{i=1}^{s}\psi^{e_i}(c_i(t-\alpha_i(\w)))\right|
&\le  \sum_{\w\in W_{\bad}}\left|\sum_{t\in\mathbb{F}_p}\prod_{i=1}^{s}\psi^{e_i}(c_i(t-\alpha_i(\w)))\right|\notag\\
&\le \sum_{\w\in W_{\bad}}p\notag\\
&=\#W_{\bad}\cdot p \notag\\
&\le \binom{s}{2}\cdot p^{m-1}.
\end{align}
		
Combining (\ref{Eq. estimate for good w}) and (\ref{Eq. estimate for bad w}) with (\ref{Eq. sum is equal to good plus bad}), we obtain 
\begin{align*}
	 &\left|\sum_{\x\in\mathbb{F}_p^m}\prod_{i=1}^{s}\chi_i(L_i(\x))\right|\\
\le &\left|\sum_{\w\in W_{\good}}\sum_{t\in\mathbb{F}_p}\prod_{i=1}^{s}\psi^{e_i}(c_i(t-\alpha_i(\w)))\right|+\left|\sum_{\w\in W_{\bad}}\sum_{t\in\mathbb{F}_p}\prod_{i=1}^{s}\psi^{e_i}(c_i(t-\alpha_i(\w)))\right|\\
\le & (s-1)\cdot p^{m-\frac{1}{2}}+\frac{s(s-1)}{2}\cdot p^{m-1}.
\end{align*}
		
In view of the above, we have completed the proof. 
	\end{proof}

\begin{thebibliography}{99}
\bibitem[A]{A} N. Alon, Combinatorial Nullstellensatz, Comb. Probab. Comput. 8 (1999), 7--29.
\bibitem[CFZ]{CFZ} D. Conlon, J. Fox and Y. Zhao, A relative Szemer\'edi theorem, Geom. Funct. Anal. 25 (2015), 733--762.
\bibitem[Cohen15]{Cohen15} S. D. Cohen, T. Oliveira e Silva and T. Trudgian, On consecutive primitive elements in a finite field, Bull. Lond. Math. Soc. 47 (2015), 418--426.
\bibitem[Cohen21]{Cohen21} S. D. Cohen, H. Sharma and R. Sharma, Primitive values of rational functions at primitive elements of a finite field, J. Number Theory 219 (2021), 237--246.
\bibitem[LN]{LN} R. Lidl and H. Niederreiter, Finite Fields, 2nd edition, Cambridge University Press, 1997.
\end{thebibliography}
\end{document}
